\section*{Author Section 2.2: quantum information}
The $\ell_1$ distance of states is $\|\rho-\sigma\|_1=\operatorname{Tr}\sqrt{(\rho-\sigma)^\dagger(\rho-\sigma)}=\operatorname{Tr}|\rho-\sigma|$. Their fidelity is $\mathrm F(\rho,\sigma)=\|\sqrt\rho\sqrt\sigma\|_1$.
\paragraph{Fact 13 (Fuchs--van de Graaf).} $2(1-\mathrm F(\rho,\sigma))\le\|\rho-\sigma\|_1\le2\sqrt{1-\mathrm F(\rho,\sigma)^2}$. For pure states, the authors state
\[\|\proj\psi-\proj\phi\|_1=\sqrt{1-\mathrm F(\proj\psi,\proj\phi)^2}=\sqrt{1-|\langle\psi|\phi\rangle|^2}.\]
\paragraph{Fact 14 (Uhlmann).} If states on $X$ are purified to $\ket\rho,\ket\sigma$ on $XY$, then $\mathrm F(\rho,\sigma)=\max_U|\bra\rho\Id X\otimes U\ket\sigma|$, over unitaries on $Y$.
\paragraph{Fact 15 (Data processing).} For a quantum channel $\mathcal E$, $\|\mathcal E(\rho)-\mathcal E(\sigma)\|_1\le\|\rho-\sigma\|_1$ and $\mathrm F(\mathcal E(\rho),\mathcal E(\sigma))\ge\mathrm F(\rho,\sigma)$.
Entropy is $S(\rho)=-\operatorname{Tr}(\rho\log\rho)$. Relative entropy is $S(\rho\|\sigma)=\operatorname{Tr}(\rho\log\rho)-\operatorname{Tr}(\rho\log\sigma)$. Relative min-entropy is $S_\infty(\rho\|\sigma)=\min\{\lambda:\rho\le2^\lambda\sigma\}$. These are finite only if the support of $\rho$ is contained in that of $\sigma$, and $0\le S(\rho\|\sigma)\le S_\infty(\rho\|\sigma)$. The authors define the smooth relative min-entropy by
\[S_\infty^\epsilon(\rho\|\sigma)=\inf_{\rho':\|\rho-\rho'\|_1\le\epsilon}S(\rho'\|\sigma).\]
It can be finite even if the original supports are not contained. Although $S_\infty$ cannot be upper-bounded by $S$, its smoothed version can, by the quantum substate theorem.
\paragraph{Fact 16 (Quantum Substate Theorem [31,23]).} For states with $\operatorname{supp}(\rho)\subseteq\operatorname{supp}(\sigma)$ and $\epsilon>0$,
\[S_\infty^\epsilon(\rho\|\sigma)\le\frac{4S(\rho\|\sigma)}{\epsilon^2}+\log\frac1{1-\epsilon^2/4}.\]
\paragraph{Fact 17 (Pinsker).} $\|\rho-\sigma\|_1\le\sqrt{S(\rho\|\sigma)}$.
\paragraph{Fact 18.} If $\sigma=\epsilon\rho+(1-\epsilon)\rho'$, then $S_\infty(\rho\|\sigma)\le\log(1/\epsilon)$.
\paragraph{Fact 19.} If $\operatorname{supp}(\rho)\subseteq\operatorname{supp}(\phi)\subseteq\operatorname{supp}(\sigma)$, then $S_\infty(\rho\|\sigma)\le S_\infty(\rho\|\phi)+S_\infty(\phi\|\sigma)$.
\paragraph{Fact 20.} For unitary $U$, $S_\infty(U\rho U^\dagger\|U\sigma U^\dagger)=S_\infty(\rho\|\sigma)$.
A CQ state is $\rho_{XY}=\sum_xP_X(x)\proj{x}_X\otimes\rho_{Y|x}$ with classical register $X$ and quantum $Y$. $X$ also denotes the associated random variable. $\rho_{Y|x}$ is the state conditioned on observing $x$ on $X$, so $\rho_{XY|x}=\proj{x}_X\otimes\rho_{Y|x}$. If registers are clear, write $\rho_x$. Mutual information is $I(Y:Z)_\rho=S(\rho_{YZ}\|\rho_Y\otimes\rho_Z)$. Conditioned on classical $X$, $I(Y:Z|X)=\mathbb E_{P_X}[I(Y:Z)_{\rho_x}]$. The chain rule is $I(XY:Z)_\rho=I(X:Z)_\rho+I(Y:Z|X)_\rho$.
\paragraph{Fact 21 ([6],LemmaB.7).} $\inf_{\sigma_Z}S_\infty(\rho_{YZ}\|\rho_Y\otimes\sigma_Z)\le2\min\{\log|\mathcal Y|,\log|\mathcal Z|\}$.
\paragraph{Fact 22.} For CQ states $\rho_{XY}=\sum_xP_X(x)\proj{x}_X\otimes\rho_{Y|x}$ and $\sigma_{XY}=\sum_xP_{X'}(x)\proj{x}_X\otimes\sigma_{Y|x}$,
\[S(\rho_{XY}\|\sigma_{XY})=S(P_X\|P_{X'})+\mathbb E_{P_X}[S(\rho_{Y|x}\|\sigma_{Y|x})].\]
\paragraph{Fact 23.} Let $\sigma_{XYZ}=\sum_{x,y}P_{XY}(x,y)\proj{x,y}\otimes\sigma_{Z|xy}$ and $\rho_{XYZ}=\sum_{x,y}P_{X'Y'}(x,y)\proj{x,y}\otimes\sigma_{Z|xy}$, with $\|P_{XY}-P_{X'Y'}\|_1\le\delta$ and $I(Y:Z|X)_\sigma\le c$. For $0<\epsilon<1/4$,
\[P_{X'Y'}\left(S_\infty^\epsilon(\sigma_{Z|xy}\|\sigma_{Z|x})>\frac{4c+1}{\epsilon^3}\right)\le\epsilon+\delta/2.\]
\begin{proof}$\mathbb E_{P_{XY}}[S(\sigma_{Z|xy}\|\sigma_{Z|x})]=I(Y:Z|X)_\sigma\le c$. Markov gives $P_{XY}(S(\sigma_{Z|xy}\|\sigma_{Z|x})>c/\epsilon)\le\epsilon$. The substate theorem implies
\begin{align*}&P_{XY}\left(S_\infty^\epsilon(\sigma_{Z|xy}\|\sigma_{Z|x})>\frac{4c+1}{\epsilon^3}\right)\\&\le P_{XY}\left(S_\infty^\epsilon(\sigma_{Z|xy}\|\sigma_{Z|x})>\frac{4c}{\epsilon^3}+\log\frac1{1-\epsilon^2/4}\right)\le\epsilon.\end{align*}
The distance assumption transfers this to $P_{X'Y'}$ with additive $\delta/2$.\end{proof}
\paragraph{Fact 24 (Quantum Raz's Lemma [3]).} For CQ $\rho_{XY},\sigma_{XY}$, with classical $X=X_1\cdots X_k$ and $\sigma$ product across all registers, $\sum_{i=1}^kI(X_i:Y)_\rho\le S(\rho_{XY}\|\sigma_{XY})$.
\paragraph{Fact 25 ([29],Lemma2).} Suppose
\[\ket\sigma_{X\widetilde XAB}=\sum_x\sqrt{P_X(x)}\ket{xx}_{X\widetilde X}\ket\sigma_{AB|x}\]
satisfies $P_X(S_\infty^\epsilon(\sigma_{B|x}\|\sigma_B)>c)\le\delta$ for $\delta>0$. There are measurement operators $\{\Pi_x\}_x$ acting only on $X\widetilde XA$ such that: (i) each succeeds with probability $\alpha=2^{-c/\delta}$ on this state, namely $\|(\Pi_x\otimes\Id B)\ket\sigma\|_2^2=2^{-c/\delta}$; (ii) $(\Pi_x\otimes\Id B)\proj\sigma(\Pi_x\otimes\Id B)$ is $\proj{xx}\otimes\rho_x$ for a state on $AB$, and
\begin{multline*}\mathbb E_{P_X}\left[\left\|\frac1\alpha(\Pi_x\otimes\Id B)\proj\sigma_{X\widetilde XAB}(\Pi_x\otimes\Id B)\right.\right.\\\left.\left.-\proj{xx}_{X\widetilde X}\otimes\proj\sigma_{AB|x}\right\|_1\right]\le\epsilon+2\delta.\end{multline*}
The authors' version differs slightly from [29]: the original uses $I(X:B)_\sigma$ and no $\widetilde X$. In that proof the mutual information is converted to $P_X(S_\infty^\epsilon(\sigma_{B|x}\|\sigma_B)>c)$ anyway, so the first change makes no difference. The same projector taking the state without $\widetilde X$ to $\proj x\otimes\proj\sigma_{AB|x}$ also takes the state with $\widetilde X$ to $\proj{xx}\otimes\proj\sigma_{AB|x}$, so the second makes no difference.
